% Part: lambda-calculus
% Chapter: introduction
% Section: computable-lambda

\documentclass[../../../include/open-logic-section]{subfiles}

\begin{document}

\olfileid{lam}{int}{lrp}
\olsection{గణనీయ ప్రమేయాలు \usetoken{S}{lambda definable}}

\begin{thm}
\ollabel{thm:computable-lambda}
ప్రతి గణనీయ పాక్షిక ప్రమేయం \tetoken{లాంబ్డాతో నిర్వచించదగిన}{!!{lambda definable}}ది.
\end{thm}

\begin{proof}
ప్రతి పాక్షిక గణనీయ ప్రమేయం~$f$ లాంబ్డా పదం~$F$చే \tetoken{లాంబ్డాతో నిర్వచించబడిన}{!!{lambda defined}}దని
చూపాలి. క్లీనీ నియత రూప సిద్ధాంతం ప్రకారం, ప్రతి ఆదిమ పునరావృత్త
ప్రమేయం ఒక లాంబ్డా పదంచే \tetoken{లాంబ్డాతో నిర్వచించబడిన}{!!{lambda defined}}దని చూపి, ఆ తరువాత
\tetoken{లాంబ్డాతో నిర్వచించదగిన}{!!{lambda definable}} ప్రమేయాలు తగిన సంయుక్తాలు, అపరిమిత అన్వేషణ కింద
సంవృతమని చూపితే సరిపోతుంది. ప్రతి ఆదిమ పునరావృత్త ప్రమేయం ఒక లాంబ్డా
పదంచే \tetoken{లాంబ్డాతో నిర్వచించబడిన}{!!{lambda defined}}దని చూపడానికి, ప్రాథమిక ప్రమేయాలు
\tetoken{లాంబ్డాతో నిర్వచించదగిన}{!!{lambda definable}}వనీ, \tetoken{లాంబ్డాతో నిర్వచించదగిన}{!!{lambda definable}} పాక్షిక ప్రమేయాలు
సంయుక్తం, ఆదిమ పునరావృత్తి, అపరిమిత అన్వేషణ కింద సంవృతమనీ చూపితే
సరిపోతుంది.
\end{proof}

మిగిలిన నిరూపణను మరింత చదవదగినదిగా చేయడానికి ఎక్కువగా వాడుకలో ఉన్న
సంకేతనాన్ని ఉపయోగిస్తాం. ఉదాహరణకు, $Mxyz$కు బదులు $M(x, y, z)$ అని
రాస్తాం. ఇది సూచనాత్మకంగా ఉన్నప్పటికీ, టైపులేని లాంబ్డా కలనశాస్త్రంలోని
పదాలకు అనుబంధ స్థానసంఖ్యలు ఉండవని గుర్తుంచుకోవాలి; కాబట్టి ఒకే పదం~$M$కు
$M(x,y)$ అని రాయడం ఎంత అర్థవంతమో, $M(x,y,z,w)$ అని రాయడమూ అంతే
అర్థవంతం. అయితే ఈ సంకేతనాన్ని వాడడం ద్వారా $M$ను మూడు చరాల ప్రమేయంగా
పరిగణిస్తున్నామని సూచిస్తాం; నిర్వచనాల వెనుక ఉద్దేశాలను ఇది మరింత
స్పష్టంగా చేస్తుంది. ఇదే విధంగా, ``$M$ను $M =
\lambd[x][\lambd[y][\lambd[z][\dots]]]$ ద్వారా నిర్వచించండి'' అని
చెప్పడానికి బదులు ``$M(x,y,z) = \dots$ ద్వారా $M$ను నిర్వచించండి''
అంటాం.

\end{document}
